\begin{document}
\frontmatter
\begin{titlepage}
\centering
\vspace*{1.2cm}
{\large\scshape UNNS Research Collective\par}
\vspace{1.5cm}
{\fontsize{26}{31}\selectfont\bfseries Generative Information Geometry\\[0.25em] in the UNNS Substrate\par}
\vspace{0.7cm}
{\Large Ancestry, Admissibility, Topology\\ and Information Recoverability\par}
\vspace{1.4cm}
{\color{accent}\rule{0.45\textwidth}{0.8pt}}\par
\vspace{0.6cm}
{\large Research monograph\par}
{\large Version 2.0\par}
{\large November 2026\par}
\vfill
\begin{minipage}{0.88\textwidth}\small\raggedright
\textbf{How to cite.} UNNS Research Collective, \emph{Generative Information Geometry in the UNNS Substrate: Ancestry, Admissibility, Topology and Information Recoverability}, version 2.0, November 2026.\\[0.5em]
\textbf{Reference collection.} The arithmetic operator on which the worked realisation rests is documented in the frozen \emph{Mathematical Foundations of the UNNS Substrate} reference collection, cited at its pinned commit in the bibliography~\cite{UNNSFound}. The releases from which the monograph was built are published at a pinned commit~\cite{UNNSGIG}.\\[0.5em]
\textbf{Package.} The archive that ships with this monograph contains the \LaTeX{} source, the figure generators, the exact checkers, the theorem crosswalk, the proof audit and the page audit of this version.
\end{minipage}
\end{titlepage}
\chapter*{Abstract}\addcontentsline{toc}{chapter}{Abstract}
\markboth{Abstract}{}
In the UNNS Substrate a number is a \emph{generated occurrence}: it is the endpoint of histories that apply operators one after another, and the same endpoint is generally reached by many histories. This monograph develops the geometry and the information theory of those histories, and reports what the Substrate's generative arithmetic turns out to contain.

For the UNNS operator $T_N(m)=m(N{+}1)^2/N$ (SELECT-II) we prove that admissibility of a history is decided by the denominators of the products of the factors $(N{+}1)^2/N$ over the \emph{prefix sets} of its operators. Consequently the admissible histories of a source form a closure system under lcm of thresholds, they sit as a growing family of vertex sets in the permutohedron, and the associated \emph{history complex} grows monotonically along divisibility of the source. Divisibility therefore acts on topology: an explicit two-dimensional cycle of the history complex persists exactly when a monotone Boolean condition on prefix sets fails, with four minimal filling families, and for the anchor contents this reads as an explicit gcd condition on the source. For every cycle of the permutohedron survival is likewise decided by that cycle's own minimal filling families; the four families and the gcd condition are particular to the anchors. A chain-level calculus lifts holes from $k$ to $k{+}1$ operators, yields certified three- and four-dimensional arithmetic holes, and characterises exactly when a lifted class depends on the chosen filling.

The information carried by a history is organised in layers that these structures separate. The ambiguity $\log_2|F_t|$ of an endpoint decomposes exactly into source, content, swap-class and within-class parts; an exact radius-two identity describes everything local two-step statistics of the swap graph can see; and systems with equal history counts, equal entropies and equal local statistics out to radius two can differ in their Betti numbers. What is \emph{forbidden} governs what is recoverable: by Alexander duality $\beta_2(K_A)=\dim H_1(Y_A;\F)$, and for prefix-defined systems the forbidden complex $Y_A$ is homotopy equivalent to the nerve of the dead prefix sets, so a hole in the admissible histories is a cycle among the dead sets, decided exactly by a strong-collapse core, while no count-level statistic of the nerve decides it. Finally the topology is controlled by the \emph{locality} of the constraints, not by integer divisibility: for six events, constraints on consecutive pairs never create a two-dimensional hole, whereas constraints on consecutive triples can create five independent ones.

\bigskip\noindent\textbf{Thesis.} Generative information in the UNNS Substrate is not exhausted by the multiplicity of histories. It is organised by the arithmetic of admissibility, by the topology of the admissible region and by the structure of the forbidden histories, and these three determine which properties of generation an observation can recover.
\begingroup\small\tableofcontents\endgroup
\mainmatter
